\subsection{Метрики кількісної оцінки ефективності масштабування}
Для загального аналізу похибки, що виникає при апроксимації цільового відмасштабованого зображення $I^{res} = [f(x_i^n, x_j^m)]_{i,j}$ вихідним зображенням $I^{out} = [\tilde{f}(x_h^N, x_k^M)]_{h,k}$, згенерованим запропонованим алгоритмом, можна звернутися до класичної літератури з оцінки похибок поліноміальної інтерполяції. Проте в контексті даної кваліфікаційної роботи головним завданням є вимірювання результативності розробленого методу за допомогою стандартних метрик, що традиційно використовуються в задачах обробки цифрових зображень. Зокрема, базовим критерієм оцінки є середньоквадратична похибка, яка визначається наступним чином:

\begin{equation*}
	\text{MSE}(I^{res}, I^{out}) = 
	\begin{cases} 
		\frac{1}{MN} \|I^{res} - I^{out}\|_F^2, & \text{для напівтонових}, \\[10pt]
		\begin{aligned}
			\frac{1}{3NM} \Big\{ &\|I_R^{res} - I_R^{out}\|_F^2 + \|I_G^{res} - I_G^{out}\|_F^2 + \\
			& + \|I_B^{res} - I_B^{out}\|_F^2 \Big\},
		\end{aligned} & \text{для кольорових (RGB)}, 
	\end{cases}
\end{equation*}

де $\|\cdot\|_F$ позначає норму Фробеніуса:

\begin{equation*}
	\|A\|_F := \left( \sum_{k=1}^N \sum_{h=1}^M A_{k,h}^2 \right)^{\frac{1}{2}}, \quad A = (A_{k,h}) \in \mathbb{R}^{N \times M}.
\end{equation*}

Наступною важливою метрикою, що розраховується на базі MSE, є пікове відношення корисних даних до шуму (у класичній термінології --- Peak Signal-to-Noise Ratio, PSNR):

\begin{equation}\tag{25}
	\text{PSNR}(I^{res}, I^{out}) = 20 \log_{10} \left( \frac{\max_f}{\sqrt{\text{MSE}(I^{res}, I^{out})}} \right),
\end{equation}

де параметр $\max_f = 255$ відповідає максимальному значенню яскравості пікселя за умови використання стандартного 8-бітного кодування.

Нарешті, для комплексної оцінки якості збереження контурів та структурної цілісності під час експериментів застосовується індекс структурної подібності. Для зображень у градаціях сірого ця метрика обчислюється за формулою:

\begin{equation}\tag{26}
	\text{SSIM}(I^{res}, I^{out}) = \frac{[2\mu(I^{res})\mu(I^{out}) + c_1][2\text{cov}(I^{res}, I^{out}) + c_2]}{[\mu^2(I^{res}) + \mu^2(I^{out}) + c_1][\sigma^2(I^{res}) + \sigma^2(I^{out}) + c_2]},
\end{equation}

де $\mu(A)$, $\sigma^2(A)$ та $\text{cov}(A,B)$ позначають відповідно середнє значення, дисперсію та коваріацію матриць $A$ та $B$, а $c_1, c_2$ --- стабілізуючі константи, які зазвичай обираються як $c_1 = (0.01 \times L)$ та $c_2 = (0.03 \times L)$, де $L = 255$ --- динамічний діапазон значень пікселів для 8-бітних зображень. 

У випадку кольорових RGB-зображень, спочатку здійснюється їх перетворення у колірний простір YCbCr, після чого індекс SSIM обчислюється за наведеною вище формулою виключно для каналу інтенсивності Y (яскравості).